% ============================================================
% LOGICAL DEPENDENCIES
% Viability Philosophy System
% Book I: Ontology of Viability
% ============================================================

\cleardoublepage

\section*{Logical Dependencies}

\phantomsection
\addcontentsline{toc}{section}{Logical Dependencies}

% ------------------------------------------------------------
% 1. PURPOSE AND SCOPE
% ------------------------------------------------------------

\subsection*{1. Purpose and Scope}

Logical dependencies identify and organize relations
of dependence between logical entities within the
Viability Philosophy System.

Their purpose is to make explicit how the identity,
definition, formulation, validity, or theoretical
position of one entity may depend on another entity.

The identification of a logical entity establishes
its identity within the system. The identification
of a logical dependency establishes a relation
between entities whose logical organization cannot
be adequately represented by identifying the
entities separately.

Logical dependencies therefore concern the internal
organization of the logical system rather than merely
the order in which its components appear in the text.

A dependency relation must be distinguished from
a textual reference, a thematic association, or
the mere co-occurrence of two entities in the same
discussion.

The logical dependency system serves the following
purposes:

\begin{itemize}
    \item identifying dependencies between logical
    entities;
    \item making relevant relations between entities
    explicit and traceable;
    \item distinguishing dependent entities from
    entities that are merely associated;
    \item clarifying the logical organization of
    the system;
    \item supporting the examination of the
    consequences of revising or removing an entity;
    \item coordinating the logical inventory with
    the argumentative and structural organization
    of the books.
\end{itemize}

This document establishes the general principles
for identifying, describing, recording, and
maintaining logical dependencies.

% ------------------------------------------------------------
% 2. LOGICAL ENTITIES AND THEIR RELATIONS
% ------------------------------------------------------------

\subsection*{2. Logical Entities and Their Relations}

Logical entities are identified independently
of the relations in which they participate.

An entity identifier establishes which entity
is being discussed. It does not, by itself,
specify how that entity relates to other entities.

For example:

\begin{center}
\begin{tabular}{ll}
    \toprule
    \textbf{Identifier} & \textbf{Function} \\
    \midrule
    \texttt{AX-001}
        & Identifies a logical entity \\
    \texttt{DF-001}
        & Identifies another logical entity \\
    \bottomrule
\end{tabular}
\end{center}

The existence of these two identifiers does not
establish any particular relation between the
corresponding entities.

A dependency must be separately identified
and justified.

The logical organization of the system therefore
requires both:

\begin{enumerate}
    \item the identification of its logical entities;
    \item the identification of relevant relations
    of dependence between those entities.
\end{enumerate}

The first establishes the entities that belong
to the logical inventory.

The second establishes how particular entities
are related through dependence.

% ------------------------------------------------------------
% 3. THE MEANING OF LOGICAL DEPENDENCY
% ------------------------------------------------------------

\subsection*{3. The Meaning of Logical Dependency}

A logical dependency exists when the status
or formulation of one logical entity depends
on another in a specified and substantively
relevant sense.

The nature of the dependence must be established
by examining the actual relation between the
entities.

A dependency must not be inferred solely from:

\begin{itemize}
    \item the order in which entities are introduced;
    \item the order in which they appear in a chapter;
    \item their occurrence in the same paragraph;
    \item similarity of terminology;
    \item thematic association;
    \item the fact that one entity is discussed
    after another.
\end{itemize}

These circumstances may indicate that two
entities are related, but they do not establish
a logical dependency.

A dependency requires a substantive account
of what depends on what and in what respect.

The system must therefore distinguish
between the existence of a relation and
the specific nature of that relation.

% ------------------------------------------------------------
% 4. DIRECTION OF DEPENDENCY
% ------------------------------------------------------------

\subsection*{4. Direction of Dependency}

A logical dependency must identify the
direction in which the dependence holds.

Where one entity depends on another, the
record must make clear which entity is
dependent and which entity functions
as the entity upon which the dependence
is established.

The two positions must not be reversed
or left ambiguous.

For purposes of description, a dependency
may be represented schematically as:

\begin{center}
    \texttt{Entity A} $\longrightarrow$
    \texttt{Entity B}
\end{center}

This notation is illustrative only.
It does not, by itself, establish a
universal convention for the direction
of the arrow.

Any formal notation adopted for the
system must explicitly state whether
the arrow points from the dependent
entity toward the entity on which
it depends, or in the opposite direction.

Until such a convention is formally
established, the relation should be
expressed in unambiguous prose.

The direction of a dependency must
be determined by the actual logical
relation rather than by the order
of presentation.

% ------------------------------------------------------------
% 5. DEPENDENCY AND OTHER FORMS
% OF RELATION
% ------------------------------------------------------------

\subsection*{5. Dependency and Other Forms of Relation}

Not every relation between logical
entities is a relation of dependence.

The system must preserve the distinction
between logical dependency and other
possible forms of association.

\subsubsection*{5.1. Thematic Association}

Two entities may concern the same
subject or belong to a common field
of discussion.

Such an association does not,
by itself, establish that either
entity depends on the other.

\subsubsection*{5.2. Conceptual Relation}

Two entities may be conceptually
related because their meanings,
scope, or applications intersect.

A conceptual relation may be relevant
to the interpretation of the system
without constituting a dependency.

\subsubsection*{5.3. Textual Reference}

One passage may refer to an entity
discussed in another passage.

This establishes a relation between
textual discussions or references,
but it does not automatically
establish a logical dependency
between the entities involved.

\subsubsection*{5.4. Argumentative Relation}

Two entities may occur within
the same argument or participate
in a common argumentative context.

Their argumentative association
must not be treated as sufficient
evidence of a dependency unless
the relevant dependence is
substantively established.

\subsubsection*{5.5. Logical Dependency}

A logical dependency requires
a more specific account.

It must identify the entities
involved, the direction of the
dependence, and the respect
in which the dependence holds.

The nature of the relation must
be supported by the logical
organization of the system.

% ------------------------------------------------------------
% 6. IDENTIFICATION OF DEPENDENCIES
% ------------------------------------------------------------

\subsection*{6. Identification of Dependencies}

A dependency should be recorded
only when the relation between
the relevant entities has been
sufficiently established.

The identification process must
avoid converting every conceptual
or textual association into
a dependency.

At a minimum, the identification
of a dependency requires:

\begin{enumerate}
    \item identification of the
    entities involved;
    \item determination of the
    direction of dependence;
    \item specification of the
    respect in which the dependence
    holds;
    \item an adequate basis for
    recognizing the relation
    as a logical dependency.
\end{enumerate}

Where any of these elements
remains unresolved, the relation
must not be presented as though
its logical status were already
settled.

An uncertain or proposed relation
must be distinguished from an
established dependency.

% ------------------------------------------------------------
% 7. DEPENDENCY RECORDS
% ------------------------------------------------------------

\subsection*{7. Dependency Records}

Logical dependencies should
be maintained in a form that
allows the relations between
entities to be inspected and
traced.

A dependency record should
identify the following information:

\begin{itemize}
    \item the identifier of the
    dependent entity;
    \item the identifier of the
    entity upon which it depends;
    \item the direction of the
    dependency;
    \item a description of the
    respect in which the dependency
    holds;
    \item the basis on which
    the dependency is established;
    \item the relevant textual
    locations, where these are
    necessary for verification.
\end{itemize}

The description of a dependency
must be sufficiently precise
to distinguish it from other
relations involving the same
entities.

A record must not rely exclusively
on the names of the entities
or on an unexplained symbolic
connection between their
identifiers.

The textual locations associated
with a dependency provide
traceability to the relevant
discussion. They do not replace
the logical justification
of the dependency itself.

% ------------------------------------------------------------
% 8. DEPENDENCY AND LOGICAL
% CLASSIFICATION
% ------------------------------------------------------------

\subsection*{8. Dependency and Logical Classification}

Logical entity identifiers
classify and identify entities
within the established logical
inventory.

Logical dependencies describe
relations between those entities.

The classification of an entity
and the dependencies in which
it participates are therefore
distinct aspects of the system.

An entity's membership in a
particular identifier family
does not, by itself, determine
which other entities it depends
on or which entities depend
on it.

Likewise, a dependency between
two entities does not alter
their identifier families
merely by virtue of that
relation.

The logical dependency system
must therefore operate in
coordination with the logical
entity identification system
without replacing or modifying
its classification rules.

% ------------------------------------------------------------
% 9. DEPENDENCY AND STRUCTURAL
% LOCATION
% ------------------------------------------------------------

\subsection*{9. Dependency and Structural Location}

Logical dependencies belong
to the logical organization
of the system.

They must be distinguished
from the structural locations
at which the relevant entities
or their dependencies are
discussed.

Two entities may be introduced
in different chapters and
nevertheless stand in a
dependency relation.

Conversely, two entities may
be discussed in the same
paragraph without depending
on one another.

A structural reference may
locate the passage in which
a dependency is established
or explained.

It does not constitute the
dependency itself.

The logical relation must
remain identifiable independently
of the particular location
in which it is presented.

% ------------------------------------------------------------
% 10. DEPENDENCY AND ENTITY
% LOCATION REFERENCES
% ------------------------------------------------------------

\subsection*{10. Dependency and Entity
Location References}

Entity location references
connect logical entities
with their textual occurrences.

Logical dependency records
connect logical entities
through relations of dependence.

These functions must remain
distinct.

A location record may identify
where a dependency is discussed,
but it does not establish
that the dependency exists.

Similarly, a dependency record
may identify two related
entities without specifying
every passage in which those
entities occur.

The two systems should be
coordinated so that logical
relations can be traced
to their relevant textual
treatments.

% ------------------------------------------------------------
% 11. DEPENDENCY AND ARGUMENTATIVE
STRUCTURE
% ------------------------------------------------------------

\subsection*{11. Dependency and Argumentative
Structure}

The argumentative organization
of the book may involve
premises, claims, explanations,
distinctions, and conclusions.

Some of these elements may
stand in relations of logical
dependence.

However, the structure of
an argument and the dependency
relations between logical
entities must not automatically
be treated as identical.

An argument may refer to
several entities without
establishing that each
entity depends on every
other entity.

Likewise, a logical dependency
may exist independently
of the particular argument
in which it is discussed.

Where an argument relies
on a dependency, the relevant
relation must be made explicit
rather than left to be inferred
solely from textual sequence.

% ------------------------------------------------------------
% 12. COMPLEX AND MULTIPLE
DEPENDENCIES
% ------------------------------------------------------------

\subsection*{12. Complex and Multiple Dependencies}

A logical entity may depend
on more than one other entity.

Likewise, more than one
entity may depend on the
same entity.

The existence of multiple
dependencies must be recorded
without collapsing distinct
relations into an undifferentiated
association.

Each dependency must be
identifiable in its own right.

Where an entity participates
in several dependencies,
the relevant records should
make it possible to distinguish
the entities involved and
the specific respect in which
each relation holds.

A complex pattern of dependencies
must not be simplified in
a way that conceals relevant
differences between its
component relations.

% ------------------------------------------------------------
% 13. DEPENDENCY AND LOGICAL
CONSISTENCY
% ------------------------------------------------------------

\subsection*{13. Dependency and Logical Consistency}

The identification of logical
dependencies contributes
to the examination of the
consistency and organization
of the system.

A dependency may reveal
that the formulation or
status of one entity requires
the clarification, revision,
or examination of another.

For this reason, dependencies
must be traceable and their
descriptions must remain
consistent with the actual
logical relations established
in the work.

The existence of a dependency
does not automatically establish
that the dependent entity
is valid or invalid.

It identifies a relation
whose consequences must
be examined according to
the nature of the dependence.

Any assessment of validity,
consistency, or justification
must be grounded in the
relevant logical and philosophical
analysis.

% ------------------------------------------------------------
% 14. REVISION AND DEPENDENCY
MAINTENANCE
% ------------------------------------------------------------

\subsection*{14. Revision and Dependency
Maintenance}

Changes to logical entities
may affect the dependencies
in which those entities
participate.

For this reason, the dependency
records must be reviewed
whenever an entity is
substantially revised,
reclassified, replaced,
or retired.

\subsubsection*{14.1. Revision of an Entity}

When an entity is revised
without losing its identity,
its dependency records
must be examined to determine
whether the relations remain
accurate.

A change in wording does
not automatically require
a change in the dependency
structure.

\subsubsection*{14.2. Replacement of an Entity}

When an entity is replaced
by a distinct entity, the
dependencies associated
with the earlier entity
must be reviewed.

They must not automatically
be transferred to the
replacement entity.

Each relation must be
reassessed according to
the actual logical status
of the new entity.

\subsubsection*{14.3. Removal or Retirement}

When an entity is removed
or retired, all dependencies
involving that entity must
be reviewed.

The records must reflect
whether the relevant
dependencies remain meaningful,
require revision, or no
longer belong to the
current logical organization
of the system.

\subsubsection*{14.4. Changes in Dependency}

A dependency may itself
require revision even when
the identities of the
entities involved remain
unchanged.

Such a change must be
recorded as a change in
the relation, not automatically
as a change in the identity
of either entity.

% ------------------------------------------------------------
% 15. UNRESOLVED DEPENDENCIES
% ------------------------------------------------------------

\subsection*{15. Unresolved Dependencies}

Not every proposed relation
between logical entities
will be sufficiently established
at the time it is first
considered.

Some relations may remain
under examination.

Such relations must be
distinguished from dependencies
that have been established
within the logical organization
of the system.

An unresolved relation
must not be recorded
or cited as though its
status were settled.

Where the system maintains
records of proposed or
unresolved dependencies,
their provisional status
must be made explicit.

The resolution of such
relations requires examination
of the relevant conceptual
and logical grounds.

% ------------------------------------------------------------
% 16. EXTERNAL SOURCES
% ------------------------------------------------------------

\subsection*{16. External Sources}

Logical dependencies concern
relations between entities
within the logical organization
of the Viability Philosophy
System.

External books, articles,
authors, and other sources
must not be assigned internal
logical entity identifiers
merely because they are
cited in support of, or
in connection with, an
internal discussion.

An external source may
provide evidence, criticism,
an interpretation, or
a conceptual point of
comparison relevant to
a dependency.

However, the bibliographic
identity of the source
must remain distinct from
the internal logical entities
and the dependencies
established between them.

References to external
sources are governed by
the corresponding external
source and bibliographic
referencing documents.

% ------------------------------------------------------------
% 17. RELATION TO THE OTHER
% REFERENCING DOCUMENTS
% ------------------------------------------------------------

\subsection*{17. Relation to the Other
Referencing Documents}

The logical dependency system
forms one component of
the broader referencing
architecture of the Viability
Philosophy System.

Its operation must remain
consistent with the other
referencing components.

In particular:

\begin{itemize}
    \item \textit{Structural Identifiers}
    identify textual locations
    within the hierarchy
    of a book.

    \item \textit{Logical Entity
    Identifiers} identify
    entities within the
    logical organization
    of the system.

    \item \textit{Internal
    Cross-References} connect
    textual locations and
    logical entities for
    purposes of navigation
    and reference.

    \item \textit{Entity Location
    References} record the
    textual locations at
    which logical entities
    occur.

    \item \textit{External Source
    References} govern references
    to sources outside
    the system.

    \item \textit{Bibliographic
    Organization} governs
    the organization of
    bibliographic information.
\end{itemize}

These components must be
coordinated while preserving
their distinct functions.

% ------------------------------------------------------------
% 18. SUMMARY
% ------------------------------------------------------------

\subsection*{18. Summary}

Logical dependencies identify
relations of dependence
between entities within
the Viability Philosophy
System.

They must be distinguished
from textual references,
thematic associations,
conceptual relations,
and mere co-occurrences
within the manuscript.

Every dependency must
identify the entities
involved, establish the
direction of dependence,
and specify the respect
in which the dependence
holds.

Dependency records must
remain coordinated with
the logical entity inventory
and the relevant textual
locations.

Changes to entities or
to the relations between
them require the review
of affected dependency
records.

The logical dependency
system therefore makes
the internal organization
of relations between
logical entities traceable
without confusing those
relations with textual
location, entity identity,
or bibliographic reference.

\cleardoublepage

% ============================================================
% END OF LOGICAL DEPENDENCIES
% ============================================================